% from ega3-1-fr.tex
\oldpage[III]{155}

anneau adique noethérien dont $\mathfrak K_1$ est un idéal de définition. Posons
$X_1=\operatorname{Spec}(B_1)$ et soit $h:X_1\to X$ le morphisme correspondant à
l'homomorphisme canonique $B\to B_1$~; si $X'_1=h^{-1}(X')$, on a donc
$X'_1=\mathrm V(\mathfrak K_1)$. Posons enfin
$Z_1=Z\times_X X_1=Z_{(X_1)}$, $f_1=f_{(X_1)}:Z_1\to X_1$, qui est un morphisme
propre (II, 5.4.2), et désignons par $\widehat X_1$ le complété de $X_1$ le long
de $X'_1$, par
$\widehat Z_1=Z_1\times_{X_1}\widehat X_1$ le complété de $Z_1$ le long de
$Z'_1=f_1^{-1}(X'_1)$, par $\widehat f_1$ le prolongement de $f_1$ aux complétés.
Il est immédiat que le prolongement $\widehat h:\widehat X_1\to\widehat X$ de $h$
aux complétés est un isomorphisme, correspondant à l'application identique de $B_1$
(I, 10.9.1)~; on en conclut que l'homomorphisme correspondant
$\widehat Z_1\to\widehat Z$ est aussi un isomorphisme, ces isomorphismes identifiant
$\widehat f_1$ et $\widehat f$. Enfin,
$\mathcal M_1=h^*(\mathcal M)$ est un Idéal cohérent de $\mathcal O_{X_1}$ et
\[
\operatorname{Supp}(\mathcal O_{X_1}/\mathcal M_1)
=h^{-1}(\operatorname{Supp}(\mathcal O_X/\mathcal M))
\]
(I, 9.1.13), donc, si
$U_1=X_1-\operatorname{Supp}(\mathcal O_{X_1}/\mathcal M_1)$, on a
$U_1=h^{-1}(U)$, d'où résulte aussitôt que la restriction
$f_1^{-1}(U_1)\to U_1$ de $f_1$ est un isomorphisme (I, 3.2.7)~; en outre, les
complétés $\widehat{\mathcal M}$ et $\widehat{\mathcal M}_1$ s'identifient par
$\widehat h$ (I, 10.9.5). Toutes les hypothèses de
(\hyperref[III.5.3.4-fr]{5.3.4}) sont donc remplies par $X_1$, $X'_1$, $f_1$ et
$\mathcal M_1$, et on peut donc désormais supposer $B$ adique noethérien et
$\mathfrak K$ un idéal de définition de $B$. On a alors
$\widehat X=\operatorname{Spf}(B)$, et $\mathcal F=N^\Delta$, où $N$ est un
$B$-module de type fini, d'où $\mathcal F=\widehat{\mathcal G}$, où $\mathcal G$
est le $\mathcal O_X$-Module cohérent $\widetilde N$ (I, 10.10.5), et par suite
\[
\widehat f^{,*}(\mathcal F)=\widehat{f^*(\mathcal G)}
\]
(I, 10.9.5). En outre, en vertu du premier théorème de comparaison (4.1.5),
\[
\widehat f_*\bigl(\widehat{f^*(\mathcal G)}\bigr)
\simeq\widehat{f_*\bigl(f^*(\mathcal G)\bigr)},
\]
et l'homomorphisme canonique $\rho_{\mathcal F}$ n'est autre que
$\widehat{\rho}_{\mathcal G}$ en vertu de
(\hyperref[III.5.1.3-fr]{5.1.3}). Or, le noyau $\mathcal P$ et le conoyau
$\mathcal R$ de
\[
\rho_{\mathcal G}:\mathcal G\longrightarrow f_*\bigl(f^*(\mathcal G)\bigr)
\]
sont des $\mathcal O_X$-Modules cohérents, et par hypothèse leurs restrictions à
$U$ sont évidemment nulles. Il existe par suite un entier $n>0$ tel que
$\mathcal M^n\mathcal P=\mathcal M^n\mathcal R=0$ (I, 9.3.4)~; on en conclut que
$\widehat{\mathcal M}^{,n}\widehat{\mathcal P}=
\widehat{\mathcal M}^{,n}\widehat{\mathcal R}=0$
(I, 10.8.8 et 10.8.10).
\end{proof}

\begin{env}[5.3.5]
\label{III.5.3.5-fr}
\emph{Fin de la démonstration du théorème d'existence.}
Les hypothèses étant celles de (\hyperref[III.5.1.4-fr]{5.1.4}), nous allons
utiliser le principe de récurrence noethérienne ($0_{\mathrm I}$, 2.2.2) en supposant
donc le théorème vrai pour tout sous-préschéma fermé $T$ de $X$ dont l'espace
sous-jacent est distinct de $X$ (le complété $\widehat T$ étant bien entendu le
complété de $T$ le long de $T\cap X'$). On peut supposer $X$ non vide. Comme $f$
est séparé et de type fini, on peut appliquer le lemme de Chow (II, 5.6.1)~: il
existe donc un $Y$-schéma $Z$ et un $Y$-morphisme $g:Z\to X$ tels que le morphisme
structural $h:Z\to Y$ soit quasi-projectif, le morphisme $g$ projectif et surjectif,
et en outre un ouvert non vide $U$ de $X$ tel que la restriction
$g^{-1}(U)\to U$ soit un isomorphisme. Soit $\mathcal M$ un Idéal cohérent de
$\mathcal O_X$ définissant un sous-préschéma fermé d'espace sous-jacent $X-U$
(I, 5.2.2), et soit $\mathcal F$ un $\mathcal O_{\widehat X}$-Module cohérent dont
le support $E$ est propre sur $Y$~; désignons par $\widehat Z$ le complété de $Z$
le long de $h^{-1}(Y')$, par $\widehat g:\widehat Z\to\widehat X$ le prolongement
de $g$ aux complétés. Alors $\widehat g^{,*}(\mathcal F)$ est un
$\mathcal O_{\widehat Z}$-Module cohérent dont le support est contenu dans
$g^{-1}(E)$ et est par suite propre sur $Y$, puisque $g$ est projectif, donc propre
(II, 5.4.6). Comme $h$ est quasi-projectif, $\widehat g^{,*}(\mathcal F)$ est
algébrisable en vertu de (\hyperref[III.5.2.6-fr]{5.2.6}). On en conclut

